\section{Financial reporting framework}

\begin{qfmethod}[A purpose-driven financial statement analysis process]\label{meth:statement-analysis-process}
\qffield{Statement} Reporting describes a company's position and performance; analysis uses those reports to make an equity, credit or operating decision. Start by defining the question, horizon, audience and decision rule. Collect statements, notes and independent business evidence; process them into comparable statements and ratios; interpret performance, risk and uncertainty; communicate a conclusion with assumptions; then monitor facts that could change it. An income statement covers a period, a balance sheet a date, and cash-flow and equity statements reconcile changes between dates. Accrual recognition records economic activity when it occurs rather than when cash moves. Notes are part of the statements, not optional detail \citep[chap.~1, solutions~1--19, PDF pp.~173--174]{statement-analysis-workbook-2020}; \citep[IAS~1, paras.~10, 15, 25, 27 and~106]{adopted-accounting-standards-2023}.
\qffield{Practice} Meridian Components, an invented industrial distributor, asks for a revolving credit facility. A lender first sets a USD 50 million minimum cash balance and a debt/EBITDA ceiling of 2, examines collections and inventory conversion, then forecasts debt service under a sales shock. An equity analyst instead asks whether normalized earnings and reinvestment justify a price. The same accounts support different decisions: a high reported ROE cannot by itself answer either question. The linked Meridian statements below provide the common numerical base; isolated accounting variants in later entries are explicitly identified.
\qffield{Pitfalls} Calculating ratios before identifying the decision; mixing annual flows with an unrelated balance-sheet date; confusing an accounting profit with cash available for distribution. Forecasts and comparisons require interpretation after calculation.
\end{qfmethod}

\begin{qfmethod}[Using filings, notes, audit reports and external evidence]\label{meth:filings-notes-audit-evidence}
\qffield{Statement} Read the audited annual statements together with accounting policies, estimates, commitments, contingencies, segment information and subsequent events. Management commentary explains price, volume, margins, liquidity and risks, but is management's account of them. Regulatory filings also contain risk factors, legal proceedings and controls disclosures; remuneration and ownership information can appear in a proxy statement. Interim reports improve timeliness but usually have condensed disclosures and a review rather than a full annual audit. A clean opinion provides reasonable assurance about material compliance with the identified framework; it is neither a solvency guarantee nor a prediction of profitable operations \citep[chap.~1, solutions~9--19, PDF p.~174]{statement-analysis-workbook-2020}.
\qffield{Practice} To assess Meridian's receivables, reconcile the aging note to the balance sheet, compare revenue policy with delivery/acceptance terms, and inspect post-year-end collections. Compare management's explanation of inventory growth with supplier lead times, order cancellations and industry demand. An auditor's qualification for inventory records would make gross margin and collateral estimates uncertain; an adverse opinion signals pervasive misstatement, while a disclaimer indicates insufficient basis for an opinion. Critical/key audit matters identify areas needing attention, not separate adverse opinions. Supplement filings with earnings calls, competitor reports, customer/supplier evidence, regulator notices and industry statistics, respecting their dates and incentives.
\qffield{Conditions} Audit scope, filing deadlines and report terminology depend on jurisdiction. Do not infer that every interim statement is unaudited, or that audit sampling means fraud is outside the auditor's responsibilities. Retain the actual filing version and publication date; later restatements were unavailable to an earlier investor.
\end{qfmethod}

\begin{qfmethod}[Comparing reporting frameworks and tracking effective standards]\label{meth:reporting-framework-transitions}
\qffield{Statement} Identify the reporting entity, consolidated perimeter, currency, fiscal year and governing standards before comparing numbers. IFRS Accounting Standards issued by the IASB, IFRS as endorsed by a jurisdiction, and US GAAP are distinct bases. Recognition, measurement and presentation differences can affect earnings, equity and ratios without changing the underlying business. The FASB Codification is the US nongovernmental GAAP authority; Accounting Standards Updates communicate amendments. An issued standard, a proposal and a currently effective requirement have different status \citep[chap.~2, solutions~1--15, PDF pp.~175--176]{statement-analysis-workbook-2020}.
\qffield{Conditions} Examples here concern non-financial public-company reporting for years beginning in 2026, with no voluntary early adoption unless stated. IFRS paragraphs use the EU's 2023 adopted text, supplemented by the specifically identified later amendments; this is not a claim that the original 2023 regulation is a current consolidation. US rules are supported by the cited primary amendment texts and the workbook, not by an exhaustive inspection of the live Codification. Private-company alternatives and local tax rules require separate treatment.
\qffield{Practice} IFRS 18 replaces IAS 1 for annual periods beginning on or after 1 January 2027, with early application permitted: new income-statement categories/subtotals and management-defined performance-measure disclosures can change comparison bases, and its IAS 7 amendments change the indirect-method starting point and interest/dividend classification. Apply its transition requirements and restate comparative presentation as required; do not treat 2026 issuance/adoption in the EU as universal 2026 mandatory application \citep[art.~2; IFRS~18 paras.~47, 69, 117--125 and~C1--C3]{eu-2026-338-ifrs18}. US ASU 2023-09 tax disclosures already apply to calendar-2026 public companies; ASU 2025-06 internal-use software changes are generally mandatory only from calendar 2028, with early adoption permitted. Maintain a standards watch by issue date, effective period, entity scope and transition choice, then quantify the effects rather than treating all reported growth as economic growth.
\end{qfmethod}


\begin{qfexample}[Meridian Components: reconciled operating, investing and financing accounts]\label{ex:meridian-linked-statements}
\qffield{Statement} Original USD-million case, years ended 31 December. All sales are credit sales; inventory purchases are on account. The 2026 figures are a hypothetical completed reporting year, not observed results as of October 2026. The base case and its 2027 forecast use US GAAP-style cash-flow presentation: interest and tax are operating, dividends paid are financing. There are no acquisitions, foreign-exchange movements, deferred tax, OCI, new shares, noncash investing or other accounts in this base case. Debt includes a current portion of 30 in both 2025 and 2026. The hypothetical tax rate is 25 percent; current tax is paid in the same year. These simplifying assumptions make the full reconciliation visible \citep[chap.~5, solutions~12--24, PDF pp.~186--187]{statement-analysis-workbook-2020}.
\begin{center}\begin{tabular}{lrr|lrr}
Income statement & 2025 & 2026 & Balance sheet & 2025 & 2026\\\hline
Sales & 900 & 1,000 & Cash & 100 & 87.5\\
Cost of sales & 540 & 600 & Receivables & 120 & 150\\
Cash-type operating expense & 171 & 180 & Inventory & 180 & 200\\
Depreciation & 45 & 50 & Net PP\&E & 400 & 430\\
Operating profit & 144 & 170 & Total assets & 800 & 867.5\\
Interest & 22.5 & 20 & Payables & 100 & 120\\
Pretax profit & 121.5 & 150 & Accrued operating expense & 40 & 45\\
Tax expense & 30.375 & 37.5 & Debt & 250 & 220\\
Net income & 91.125 & 112.5 & Equity & 410 & 482.5
\end{tabular}\end{center}
\qffield{Practice} For 2026, customer receipts are $1{,}000-(150-120)=970$. Purchases are $600+(200-180)=620$; supplier payments are $620-(120-100)=600$. Operating-expense payments are $180-(45-40)=175$. Direct CFO is $970-600-175-20-37.5=137.5$. Indirect CFO is $112.5+50-30-20+20+5=137.5$. Investing cash flow is capital expenditure of $-80$; financing is principal repayment $-30$ and dividends $-40$, hence net cash change $-12.5$ and closing cash 87.5. PP\&E rolls from $400+80-50=430$; debt from $250-30=220$; equity from $410+112.5-40=482.5$. Both balance sheets balance. Equity is decomposed into unchanged contributed capital of 200 and retained earnings of 210 rising to 282.5; there is no OCI. Later variants do not silently accumulate into these base statements.
\end{qfexample}
